\documentclass[11pt]{article}
\usepackage[margin=0.9in]{geometry}
\usepackage{hyperref}
\usepackage{enumitem}
\usepackage{parskip}
\setlength{\parskip}{4pt}

\hypersetup{colorlinks=true, linkcolor=black, urlcolor=blue, citecolor=black}

\title{\textbf{Surrogate Aerodynamic Modeling: Predicting Transonic\\
Airfoil Coefficients via Machine Learning}}

\author{Ashish Dasu\\
\small Khoury College of Computer Sciences, Northeastern University}

\date{March 2026}

\begin{document}

\maketitle

\begin{abstract}
High-fidelity CFD simulation of airfoil aerodynamics is computationally
expensive, taking hours per configuration. This makes it impractical
for design workflows requiring thousands of evaluations. This project
proposes training machine learning regression models as fast surrogates
to predict lift ($C_l$), drag ($C_d$), and pitching moment ($C_m$)
coefficients directly from flight conditions and airfoil geometry. The
focus is the transonic regime (Mach 0.3--0.8), which is especially
challenging due to shock waves and drag divergence, and is directly
relevant to stage separation and vehicle reentry. Four regression
models are compared (polynomial regression, random forest, gradient
boosting, and a PyTorch DNN) on the TransonicSurrogate
dataset~\cite{elrefaie2024}, a publicly available collection of 1,362
CFD simulations across 8 airfoil geometries.
\end{abstract}

\section{Introduction and Motivation}

In aircraft design, accurate aerodynamic coefficient prediction
(lift, drag, and pitching moment) is critical from early geometry
sweeps through GNC. The standard approach is high-fidelity CFD, which
gives accurate results but is slow: a single simulation can take hours
on a cluster. For design optimization, this is a bottleneck since
thousands of evaluations are needed across different Mach numbers,
angles of attack, and geometries.

The transonic regime is where this problem is worst. Between Mach 0.3
and 0.8, shocks form on the airfoil surface and interact with the
boundary layer, causing nonlinear spikes in drag (drag divergence) that
simple polynomial fits cannot capture. This regime covers stage
separation, grid fin actuation, and reentry phases for modern launch
vehicles, making it both technically challenging and operationally
important.

The goal of this project is to train ML regression models that serve as
fast aerodynamic design aids, enabling rapid exploration of the
coefficient space at a fraction of the cost of direct simulation. This
has direct implications for any workflow where aerodynamic performance
must be evaluated repeatedly, from early-stage geometry sweeps to
applications where fast coefficient lookup is critical.

\section{Dataset}

This project uses the \textbf{TransonicSurrogate}
dataset~\cite{elrefaie2024}, freely available on GitHub
(\texttt{Input\_Data.csv} in the repository). The data was generated
using the \texttt{rhoCentralFoam} compressible flow solver in
OpenFOAM, solving the compressible Euler equations at roughly 7 minutes
per simulation on 128 CPU cores with 4 NVIDIA A100 GPUs.

\textbf{Dataset size:} 1,362 CFD simulations across 8 airfoil
geometries: RAE2822, RAE5212, NACA0012, NACA2412, NACA4412, NACA23012,
NACA24112, and NACA25112. These span both symmetric and cambered
profiles across a Mach range of 0.3--0.8 and angles of attack from
$-4^\circ$ to $+8^\circ$.

\textbf{Features:} Each sample includes Mach number, angle of attack,
and geometric parameters describing the airfoil shape (thickness,
camber, and related profile characteristics). The authors use a
numeric parameterization rather than one-hot encoding the airfoil
identity~\cite{elrefaie2024}.

\textbf{Targets:} Three simultaneous regression outputs per sample:
$C_l$, $C_d$, and $C_m$.

\textbf{Preprocessing:} Z-score normalization on all inputs and
per-coefficient normalization on targets. The dataset is split 60/20/20
into train, validation, and test sets per the original paper.
For generalization analysis, each trained model is additionally
evaluated on samples from a single held-out airfoil geometry (RAE2822
withheld from training) to assess performance on unseen profiles.

\section{Proposed Methods}

This is a multi-output regression problem. Four models are compared
with increasing complexity:

\begin{enumerate}[leftmargin=*, nosep]
  \item \textbf{Polynomial Regression (baseline):} Linear regression
  with degree-2 and degree-3 polynomial basis function expansion.
  Degree selected via cross-validation. Tests whether the transonic
  aerodynamic response can be captured with smooth polynomial structure.

  \item \textbf{Random Forest Regressor~\cite{breiman2001}:}
  Multi-output ensemble of decision trees (sklearn). Tuning covers
  number of trees $n \in \{50, 100, 200\}$ and max depth
  $d \in \{5, 10, \infty\}$ via 5-fold CV.

  \item \textbf{Gradient Boosting (XGBoost~\cite{chen2016}):} One
  boosted tree model per output. Hyperparameters: learning rate
  $\eta \in \{0.01, 0.05, 0.1\}$, max depth, and number of estimators,
  tuned via CV. Also provides feature importance.

  \item \textbf{Deep Neural Network (PyTorch):} Fully-connected MLP
  with a shared trunk and three output heads for $C_l$, $C_d$, $C_m$.
  Architecture: 3--5 hidden layers, width 64--256, ReLU activations,
  batch normalization, Adam optimizer, MSE loss, dropout
  ($p \in \{0.1, 0.3\}$). Final architecture chosen via validation
  loss.
\end{enumerate}

\textbf{Evaluation:} RMSE and $R^2$ per output on the held-out test
set. Predicted drag polars ($C_l$ vs.\ $C_d$) are compared against
CFD ground truth to check physical plausibility, and inference latency
is measured to quantify the speedup over direct simulation.
XGBoost and the DNN are expected to outperform polynomial regression
given the sharp nonlinearities in the transonic regime; the polynomial
baseline serves primarily to quantify exactly where smooth fits break
down.

\begin{thebibliography}{9}

\bibitem{elrefaie2024}
M.~Elrefaie, T.~Ayman, M.~Elrefaie, E.~Sayed, M.~Ayyad, and M.~M.~AbdelRahman,
``Surrogate Modeling of the Aerodynamic Performance for Airfoils in Transonic Regime,''
\textit{AIAA SciTech Forum}, 2024.
Dataset: \url{https://github.com/Mohamedelrefaie/TransonicSurrogate}

\bibitem{chen2016}
T.~Chen and C.~Guestrin,
``XGBoost: A Scalable Tree Boosting System,''
\textit{ACM SIGKDD International Conference on Knowledge Discovery
and Data Mining}, 2016.

\bibitem{breiman2001}
L.~Breiman,
``Random Forests,''
\textit{Machine Learning}, vol.~45, pp.~5--32, 2001.

\end{thebibliography}

\end{document}
